\section{Discussion}\label{sec:discussion}

\subsection{Experimental implementation}

The ablation (Section~\ref{sec:ablation}) simplifies the implementation considerably: no per-line modulators are needed. A
$200~\mathrm{GHz}$ input comb locked to the resonator FSR is split in two; one branch stays static and the other passes one intensity
or phase modulator carrying the signal. Each branch passes a phase-only (or amplitude-and-phase) spectral shaper that writes the biases
$b_\mu$ and weights $w_\mu$ of Eq.~\eqref{eq:mask}, and the recombined comb enters the resonator with the pump.
Commercial liquid-crystal-on-silicon spectral shapers set the amplitude and phase of each channel independently, with passbands down
to about $10~\mathrm{GHz}$ set in $1~\mathrm{GHz}$ steps, far below the line spacing~\cite{coherent_waveshaper}; C-band models cover
about $5~\mathrm{THz}$ (1527 to $1568~\mathrm{nm}$), enough for 16 input lines ($3.2~\mathrm{THz}$), while 32 lines or more need wider
models; their insertion loss (typically about $5~\mathrm{dB}$, up to 6.5 to $7~\mathrm{dB}$) must be included in the input power
budget. The readout is a demultiplexer after the drop port and 31 photodiodes, the same structure as frequency-multiplexed optical
reservoirs~\cite{butschek2020,lupo2023}. The hardest part is the input comb itself: it must be aligned with the resonator FSR and
phase-stable; we did not simulate its noise.

The requirement on the source follows from Section~\ref{sec:laser}: the intrinsic linewidth of the pump laser relative to the
resonance must be below a few kilohertz. This rules out ordinary DFB lasers ($100~\mathrm{kHz}$ to $1~\mathrm{MHz}$) and calls for a
narrow-linewidth laser or self-injection locking~\cite{shen2020,jin2021}. Slow frequency drift (much longer than one symbol) must be
held by a lock to within about 0.3\% of the half-linewidth; otherwise the trained readout fails as the device drifts (in the
closed-loop experiment of Appendix~\ref{app:closedloop}, a slow drift of 1\% already makes the tracking error 2.4 times worse).

\subsection{Power and rate}

The best single ring needs about $180~\mathrm{mW}$ of total optical input power (stronger drop coupling), and the best points of the
other mask variants lie between about 20 and $200~\mathrm{mW}$; a Si$_3$N$_4$ ring would need about a quarter of that, but then fewer
photons reach the detectors as well, requiring better detectors or an optical amplifier (Section~\ref{sec:platform}). The rate is
limited by the $Q$ needed for the Kerr nonlinearity to about $1~\mathrm{GBd}$. At $1~\mathrm{GBd}$ and $180~\mathrm{mW}$ the optical
energy per symbol is about $180~\mathrm{pJ}$, before the wall-plug efficiency of the laser, the 31 detection channels and the readout
computation. We do not compare energy with digital electronics, because that needs a concrete system scenario, and since the resonator
is not ahead in accuracy (Section~\ref{sec:software}), an energy benefit can only be discussed where the signal is optical and
digitisation can be avoided.

\subsection{Limitations}

\emph{Simulation only.} Our conclusions rest on the LLE and on the noise model. We chose the detector noise from datasheets
($0.3~\mu\mathrm{W}$ per line, in the middle of the 0.13 to $0.63~\mu\mathrm{W}$ of 1-GHz-class InGaAs PIN receivers), the laser
intensity noise and the frequency noise, and checked the latter against realistic spectra; thermorefractive noise of the ring, ASE beat
noise of optical amplifiers, detector quantum efficiency and the noise of the input comb are not modelled. Our implementation of the
literature schemes is also uncertain: the original papers do not give all of the cavity length, the $Q$ split and the sampling
interval, so the comparison should be read as ``in our implementation''.

\emph{Readout.} The readout quantization was computed only for the best single ring and one calibration scheme
(Section~\ref{sec:quant}); a readout weighted in the optical domain was not simulated.

\emph{Tasks.} The three benchmarks and the fibre-link task are standard; the mask's advantage concentrates on prediction that needs strong
nonlinearity and short memory, and tasks with long memory rely on external delay lines (Section~\ref{sec:memory}).

\subsection{Directions that did not work out}

Looking for an advantage of the resonator over software, we tried several more directions, all with negative or limited results. We
summarise them so that others need not repeat them:
\begin{itemize}
\item \emph{Closed-loop decisions} (Appendix~\ref{app:closedloop}): the resonator predicts the Hénon map and a controller acts on the
      prediction. The resonator works in the loop, with decision quality set by the prediction noise; it is slightly ahead of the ESN
      and NG-RC only when the software input noise is $3\times10^{-3}$, and no longer below about $2.4\times10^{-3}$.
\item \emph{A resonator controlling another Kerr resonator}: with a single-mode Kerr cavity as the plant, the resonator controller's
      tracking error ($0.023$) equals that of a 31-node ESN ($0.023$) and of a proportional--derivative controller with feedforward
      ($0.022$), and is only about 3 times better than a plain proportional--derivative controller ($0.068$); when holding the unstable
      state, the proportional--derivative controller and the ESN succeed more often (94\% and 91\%) than the resonator (68\%).
\item \emph{A shared reference ring}: a ring without signal measures the common-mode detuning jitter for subtraction in the readout.
      Even an ideal reference gains only 2\% (Section~\ref{sec:reference}).
\item \emph{Sensing}: the same comb estimates an applied detuning offset (rms error in half-linewidths). Over a narrow range ($\pm0.3$
      half-linewidths) reading 31 lines is only about 1.5 times better than a classical low-power pump-plus-reference reading
      ($1.0\times10^{-3}$ against $1.5\times10^{-3}$); over a wide range ($\pm2$ half-linewidths) the comb is far better than reading the
      pump alone ($0.07$ to $0.09$ against above $0.5$), but a low-power comb with almost no Kerr effect is better still ($0.042$), so the
      benefit comes from probing with many lines, not from the Kerr nonlinearity.
\end{itemize}
These results agree with Section~\ref{sec:software}: in accuracy, what a Kerr-microcomb reservoir can reach is close to a small
software reservoir with quantization noise; its value can only come from the physical layer (an optical signal, latency, energy),
which this paper does not demonstrate and which we leave as an open question.
